% ==========================================================
% ANEXOS
% Anexo A reproduce el instrumento de la Lección 1. Añade
% Anexo B, C, etc. como nuevas subsecciones si el curso lo pide.
% ==========================================================

\section*{\centering ANEXOS}

\subsection*{Anexo A. Instrumento de selección de idea de investigación (Lección 1)}
\notafase{Información diligenciada por el equipo en el instrumento de la Lección 1.}

\bigskip
\noindent\textbf{Datos del equipo}
\begin{center}
\begin{tabularx}{\textwidth}{L{5cm} X}
\toprule
\textbf{Campo} & \textbf{Respuesta}\\
\midrule
Integrante 1 & Gabriel Alejandro Gamba Leguizamón\\
Integrante 2 & José Fernando Barreto Franco\\
Integrante 3 & Daniel Felipe Molina Barajas\\
Fecha & 23 de agosto de 2026\\
Tema de interés general & Salud Digital, Ergonomía y Bienestar Físico en Estudiantes de Ingeniería de Sistemas\\
\bottomrule
\end{tabularx}
\end{center}

\bigskip
\noindent\textbf{Idea de investigación tentativa}

\noindent Quiero saber cuál es la relación entre el número de horas diarias
dedicadas a la exposición a pantallas, las posturas adoptadas durante las jornadas
de estudio y la frecuencia de síntomas físicos (fatiga visual y molestias
osteomusculares) en los estudiantes de Ingeniería de Sistemas de la UPTC, porque
eso importaría para proponer una guía práctica de ergonomía y pausas activas
adaptada a las rutinas de los futuros profesionales de software.

\bigskip
\noindent\textbf{Interés y pertinencia}
\begin{center}
\begin{tabularx}{\textwidth}{L{5.5cm} X}
\toprule
\textbf{Pregunta} & \textbf{Respuesta del equipo}\\
\midrule
¿Cuál es su interés personal en este tema? & Como estudiantes de Ingeniería de Sistemas, pasamos largas jornadas frente al computador y frecuentemente experimentamos fatiga visual o dolores de espalda y cuello, por lo que queremos entender cómo nuestras rutinas afectan la salud física.\\
¿Cuál es su interés académico o profesional? & El tema se relaciona directamente con la salud ocupacional e higiene postural en el ámbito de la ingeniería de software, permitiéndonos adquirir conocimientos sobre la prevención de riesgos físicos derivados del trabajo sedentario y continuo con computadores.\\
¿Es un tema prioritario o significativo para la disciplina? & Sí. La extensión de jornadas de programación y clases virtuales aumenta la incidencia de síndromes como el Síndrome Visual Informático y trastornos musculoesqueléticos en estudiantes e ingenieros, impactando su rendimiento y calidad de vida.\\
¿A quién afecta el problema que identificaron? & Principalmente a los estudiantes de la Escuela de Ingeniería de Sistemas, docentes, y profesionales del área que pasan más de 6 horas diarias expuestos a pantallas.\\
\bottomrule
\end{tabularx}
\end{center}

\bigskip
\noindent\textbf{Originalidad}
\begin{center}
\begin{tabularx}{\textwidth}{L{5.5cm} X}
\toprule
\textbf{Pregunta} & \textbf{Respuesta del equipo}\\
\midrule
¿Ya se ha investigado esto antes? & Existen estudios generales sobre ergonomía y uso de pantallas en trabajadores de oficina, pero pocos enfocados específicamente en la población universitaria de Ingeniería de Sistemas de nuestro contexto regional.\\
Si ya existe, ¿qué preguntas quedan sin responder? & Cómo se relacionan específicamente el número acumulado de horas continuas de programación con la aparición temprana de molestias osteomusculares y visuales en estudiantes jóvenes.\\
¿Qué aportarían ustedes que no exista ya? & Datos primarios y actualizados del contexto universitario local (UPTC) junto con un análisis correlacional sencillo entre hábitos posturales, horas de pantalla y síntomas físicos reportados.\\
\bottomrule
\end{tabularx}
\end{center}

\bigskip
\noindent\textbf{Viabilidad}
\begin{center}
\begin{tabularx}{\textwidth}{L{5.5cm} X}
\toprule
\textbf{Pregunta} & \textbf{Respuesta del equipo}\\
\midrule
¿Qué tan complejo es el tema para desarrollarlo en 16 semanas? & Muy viable y acotado. Consiste en diseñar y aplicar una encuesta estructurada (usando escalas validadas como la escala de usabilidad o dolor sintomático) y analizar los datos correlacionales.\\
¿Tienen acceso a los datos, herramientas o participantes que necesitarían? & Sí. Tenemos acceso directo a nuestros compañeros de carrera como participantes del estudio y utilizaremos herramientas gratuitas como Google Forms para la recolección de datos y programas como Excel para el análisis de frecuencia.\\
¿Tienen (o pueden adquirir en el curso) los conocimientos técnicos necesarios? & Sí. El diseño de encuestas y el análisis estadístico descriptivo/correlacional básico se abordan en el marco del curso de Metodología de la Investigación.\\
¿El tiempo disponible (4 h de clase + 4 h autónomas por semana) alcanza para esta idea? & Sí. El tiempo es suficiente para la revisión de antecedentes, diseño del cuestionario, recolección de respuestas en 2 semanas, análisis estadístico y redacción del informe final.\\
\bottomrule
\end{tabularx}
\end{center}

\bigskip
\noindent\textbf{Semáforo de la idea}
\begin{center}
\begin{tabularx}{\textwidth}{L{2.5cm} L{6cm} X}
\toprule
\textbf{Señal} & \textbf{Significa} & \textbf{¿Aplica a la idea?}\\
\midrule
\textcolor{greensignal}{\textbf{Verde}} & Cumple los tres criterios sin ajustes mayores. & \textbf{Sí} --- puede avanzar a la Lección 2 (antecedentes).\\
Amarillo & Cumple en general, pero un criterio necesita ajustarse. & No\\
Rojo & Falla claramente en al menos un criterio. & No\\
\bottomrule
\end{tabularx}
\end{center}
